% !TEX root = ../main.tex
\section{Introdução}
\label{sec:intro}

Cerca de metade da biomassa vegetal é carbono, fração de referência adotada
pelo \textit{Intergovernmental Panel on Climate Change} para inventários
nacionais de gases de efeito estufa (IPCC, 2019). As florestas tropicais
intactas armazenam, globalmente, parte substancial desse estoque terrestre, e
a Amazônia concentra uma das maiores reservas de carbono florestal do planeta
(PAN et al., 2011; SAATCHI et al., 2011; FAO, 2020). Quantificar a biomassa
acima do solo (\textit{aboveground biomass}, AGB) com precisão espacial é o
que sustenta inventários nacionais de emissões segundo as diretrizes
metodológicas do IPCC (IPCC, 2019), mecanismos de compensação como o REDD+,
mercados de crédito de carbono e o acompanhamento de metas do Acordo de Paris
(UNFCCC, 2015). Sem estimativas confiáveis e espacialmente distribuídas, não é
possível distinguir floresta intacta de degradada, dimensionar perdas por
desmatamento, nem validar os modelos do ciclo global de carbono que dependem
desses números como entrada (IPCC, 2021).

O gargalo histórico dessa estimativa é que AGB depende fortemente da
\textbf{estrutura vertical} da floresta --- altura das árvores, densidade e
arquitetura do dossel ---, exatamente a dimensão que sensores ópticos orbitais
(Landsat, Sentinel) enxergam mal: eles registram a reflectância do topo do
dossel e saturam em florestas densas de alta biomassa, como a maior parte da
Amazônia intacta. O LiDAR (\textit{Light Detection and Ranging}) resolve esse
problema na origem --- pulsos laser penetram o dossel e retornam um perfil
tridimensional da vegetação, medindo diretamente altura e estrutura vertical.
Isso o torna substancialmente mais acurado que índices espectrais em floresta
tropical de alta biomassa e permite calibrar relações alométricas (métricas
LiDAR $\leftrightarrow$ AGB de campo) que, em princípio, podem ser estendidas no
espaço e no tempo --- reduzindo a dependência de inventários de campo, caros e
logisticamente inviáveis na escala da bacia amazônica.

Para operacionalizar esse objetivo, a biomassa de referência desta tese é
calculada pelas mesmas equações alométricas que sustentam boa parte da
literatura de carbono florestal tropical: Chave et al. (2014) para árvores
vivas, Chambers et al. (2000) para árvores mortas em pé e Goodman et al.
(2013) para palmeiras --- a mesma base alométrica usada, por exemplo, por
Longo et al. (2016) na Amazônia.

A literatura de estimativa de AGB por LiDAR aerotransportado em floresta tropical
é madura e compartilha um padrão metodológico que vale destacar frente a este
trabalho: Asner \& Mascaro (2014) --- 804 parcelas de 1~ha em quatro países
tropicais, $R^2$ = 0,90 usando altura de topo de dossel (TCH) combinada a área
basal e densidade de madeira --- validam por \textbf{sítio}, não por parcela
embaralhada aleatoriamente. Longo et al. (2016), na Amazônia brasileira (359
parcelas, intacta e degradada), reporta $R^2 \approx$ 0,70 por regressão sobre
métricas LiDAR agregadas --- mais próximo em escopo e escala deste trabalho, mas
também sem validação k-fold aleatória explícita. Bispo et al. (2020) --- atribuída
ao Cerrado, não à Amazônia, após a correção de erro desta auditoria
(Seção~\ref{sec:lit-check}) --- reporta RMSE = 67,6~Mg/ha, rRMSE = 36,0\%,
$R^2$ = 0,58 usando só métricas LiDAR, números na mesma ordem de grandeza dos
obtidos aqui (RMSE $\approx$ 73--79~Mg/ha, rRMSE $\approx$ 43--46\%, variante sem
outliers), apesar do bioma diferente. d'Oliveira et al. (2012), em floresta
manejada no Acre, reporta $R^2$ = 0,63--0,72 em escala de \textbf{talhão}
(1.000~ha) --- suporte espacial muito maior que a parcela individual usada aqui,
o que limita a comparabilidade direta apesar da mesma região amazônica.

Mais recentemente, a literatura passou a incorporar aprendizado profundo
diretamente sobre a nuvem de pontos bruta: Oehmcke et al. (2024), em floresta
temperada dinamarquesa, encontra uma CNN de Minkowski superando PointNet, KPConv
e uma abordagem \textit{area-based} clássica --- o oposto do padrão encontrado
neste trabalho (Seção~\ref{sec:arch}), em que a abordagem parcimoniosa vence a
PointNet por uma margem grande. A leitura mais provável não é uma contradição
entre os dois resultados, mas um efeito de \textbf{regime de amostra}: florestas
temperadas manejadas tendem a ter conjuntos de treino bem maiores e menos
heterogêneos estruturalmente que os 325 pares LiDAR-inventário amazônicos
disponíveis aqui --- consistente com a leitura central deste trabalho: sob
orçamento de amostra reduzido, o efeito da complexidade representacional da
arquitetura compete com --- e pode ser dominado por --- decisões de
pré-processamento e engenharia de \textit{feature} do dado de entrada.

Os produtos orbitais da NASA (GEDI L4A/L4B) operam em escala e sensor distintos
--- LiDAR espacial \textit{footprint} ($\sim$25~m) ou grade de 1~km, não nuvem
aérea densa por parcela --- e servem aqui como baliza de ordem de grandeza, não
como comparação direta: a validação independente do GEDI L4A contra LiDAR aéreo
em savana africana (Li et al., 2024) reporta $R^2 \approx$ 0,42, abaixo de todos
os modelos desta tese na variante sem outliers, mas sobre um problema de
estimação substancialmente mais difícil (resolução espacial muito menor, amostra
espacialmente esparsa).

Um ponto, porém, permanece pouco caracterizado empiricamente para o contexto
amazônico. A literatura revisada acima mistura, na prática, dois eixos que
raramente são testados em separado: \textbf{de um lado}, o quanto a arquitetura
de modelagem em si --- de estatísticas agregadas simples a redes profundas
sobre geometria bruta --- importa para o desempenho; \textbf{de outro}, o
quanto decisões de pré-processamento e engenharia de \textit{feature} do dado
de entrada (curadoria geométrica de parcela, projeção de relevo, filtros de
altura, correção temporal) importam. Quando os dois efeitos são empilhados num
único pipeline, como é comum na literatura, não é possível atribuir o ganho de
desempenho observado a um ou a outro --- e, sob um orçamento de amostra
pequeno (nesta tese, no máximo 493 parcelas, 325 após controle de qualidade
geométrica), essa indistinção importa: um resultado que parece confirmar o
valor de uma arquitetura sofisticada pode, na verdade, ser inteiramente
explicado por uma etapa de curadoria de dados anterior ao modelo.

\label{sec:rq}
Este trabalho investiga uma pergunta específica, deliberadamente mais restrita do
que a aspiração de ``mapear a biomassa da Amazônia'':

\begin{quote}
\textbf{Num regime de poucas centenas de parcelas rotuladas, o que domina a
capacidade de um modelo generalizar na estimativa de AGB por LiDAR --- a
complexidade representacional da arquitetura ou o estágio de
pré-processamento/engenharia de \textit{feature} do dado de entrada --- e dá
para separar esses dois efeitos?}
\end{quote}

A hipótese de trabalho, motivada pela teoria estatística clássica de
viés-variância sob alta dimensionalidade e poucas amostras, é que, neste
regime de amostra reduzida, o estágio de curadoria e pré-processamento do dado
de entrada tem efeito de magnitude igual ou maior que a escolha da
arquitetura de modelagem --- e que representações de alta capacidade (PointNet
sobre nuvem crua, CNN 3D sobre voxels) sofrem desproporcionalmente mais de
sobreajuste que representações parcimoniosas (estatísticas agregadas de
altura), independentemente da riqueza estrutural que a geometria bruta, em
princípio, carrega. Testar essa hipótese exige um desenho experimental que
isole os dois eixos --- pré-processamento e arquitetura --- em vez de tratá-los
como uma correção única, como detalhado no protocolo de validação
(Seção~\ref{sec:protocolo}).

Essa pergunta se desdobra em quatro objetivos específicos, cada um respondido
numa subseção dedicada dos Resultados (Seção~\ref{sec:resultados}):

\begin{enumerate}
  \item \textbf{Eixo pré-processamento}: isolar o efeito do controle de
    qualidade geométrica das parcelas e da correção do gap temporal entre voo
    LiDAR e medição de campo sobre o desempenho dos modelos, testando cada
    tratamento isoladamente em vez de empilhá-los numa correção única
    (\textbf{§\ref{sec:preproc}}).
  \item \textbf{Eixo arquitetura}: comparar, sob o mesmo protocolo de
    validação e a mesma variante de pré-processamento, seis abordagens de
    modelagem que vão de estatísticas agregadas interpretáveis (ABA, Random
    Forest) a aprendizado profundo sobre a geometria bruta (PointNet, CNN 2D,
    CNN 3D), identificando qual representação generaliza melhor neste regime de
    amostra pequena (\textbf{§\ref{sec:arch}}).
  \item \textbf{Dentro do eixo arquitetura}, em grão mais fino: medir como o
    desempenho do GBR varia com a dimensionalidade da representação de altura
    do dossel (de 8 a 1024 atributos), para testar se mais informação bruta
    realmente se traduz em mais sinal preditivo ou se é absorvida como ruído
    de sobreajuste (\textbf{§\ref{sec:ablation}}).
  \item Identificar, através do único modelo desta comparação com importância de
    variável interpretável (ABA), quais propriedades estruturais do dossel mais
    sustentam a previsão de AGB (\textbf{§\ref{sec:interpret}}).
\end{enumerate}

Por rigor, delimita-se desde já explicitamente o que \textbf{não} está dentro do escopo do
resultado central: a validação aqui é k-fold aleatório, não
\textit{leave-one-site-out} --- portanto mede ajuste dentro da amostra de sítios
já vistos no treino, não generalização espacial para floresta amazônica não
amostrada. Há, agora, um intervalo de predição por parcela (conformal prediction
sobre resíduos out-of-fold, Seção~\ref{sec:conformal}) --- mas ele cobre só a
incerteza do modelo estrutura$\rightarrow$AGB, não o erro alométrico da fórmula de
referência nem a incerteza de medição de campo propagados até o alvo de treino.
Uma análise de sensibilidade (Seção~\ref{sec:allom-sens}) mostra, além disso, que
o próprio resultado ``CNN 2D é a melhor arquitetura'' não é robusto à escolha da
fórmula alométrica de referência --- sob densidade de madeira específica por
espécie (M3) em vez de uniforme (M1), o Random Forest assume a liderança. Essas
lacunas são discutidas com profundidade na Seção~\ref{sec:discussao} e não são
triviais: são precisamente o que separa ``a estrutura do dossel explica biomassa
nesta amostra'' de ``este pipeline está pronto para mapear a Amazônia'' --- a
segunda afirmação é mais forte do que os dados aqui sustentam.

O restante do artigo está organizado da seguinte forma. A
Seção~\ref{sec:dados} descreve a proveniência dos dados, o pipeline de
pré-processamento auditado, a biomassa de referência e as seis abordagens de
modelagem comparadas. A Seção~\ref{sec:resultados} apresenta os resultados: o
efeito do controle de qualidade geométrica e da correção de gap temporal, a
comparação entre arquiteturas, o diagnóstico do melhor modelo e a ablação de
dimensionalidade que motiva o argumento central do trabalho. A
Seção~\ref{sec:discussao} discute criticamente os resultados --- incluindo suas
limitações de validação --- e os posiciona frente à literatura e ao produto
orbital GEDI da NASA.
